\subsection{Maxima before conditional averaging}

The next local rule joins fields on actual open neighborhoods. It will be
used in both propagation arguments after their common singularities have
been canceled.

\begin{lemma}[Maximum of weak subsolutions]
\label{lem:maximum}
Let $\mathcal O\subset\R^3$ be open, let $g$ be locally bounded, and let
$F(y,t)$ be measurable in $y$, continuous in $t$, and bounded on compact
spatial sets and bounded ranges of $t$.  If finitely many functions
$v_i\in C(\mathcal O)\cap H^1_{\mathrm{loc}}(\mathcal O)$ satisfy
\[
 \Delta v_i\ge g+F(y,v_i)
\]
in distributions, their maximum satisfies the same inequality with its
own value in the reaction term.  The assertion also holds for functions
$v_i=V+a_i$ with continuous locally $H^1$ offsets $a_i$, when all
inequalities have the common singular term $-4\pi\nu$ and the reaction
vanishes in a neighborhood of $\operatorname{supp}\nu$.
\end{lemma}

\begin{proof}
It suffices to consider two functions.  For a nonnegative compactly
supported smooth test $\varphi$, choose an increasing Lipschitz function
$\theta_\eta$ that is zero on $(-\infty,0]$ and one on
$[\eta,\infty)$.  Test the first inequality with
$\theta_\eta(v_1-v_2)\varphi$ and the second with
$(1-\theta_\eta(v_1-v_2))\varphi$.  These nonnegative compactly
supported $H^1$ tests are legitimate by mollification, since the right
sides are bounded on their support.  The sum of the left sides is
\begin{align*}
 &-\int\bigl[\theta_\eta(v_1-v_2)\nabla v_1
       +(1-\theta_\eta(v_1-v_2))\nabla v_2\bigr]\cdot\nabla\varphi\\
 &\hspace{8mm}
 -\int\theta_\eta'(v_1-v_2)|\nabla(v_1-v_2)|^2\varphi.
\end{align*}
The second term is nonpositive.  Deleting it and letting $\eta\downarrow0$
gives the gradient of $\max(v_1,v_2)$ by the Sobolev positive-part
formula.  Dominated convergence gives the reaction at the maximum;
the reactions agree at ties.  Induction proves the finite version.

For the last assertion subtract $V$ in each inequality.  The singular
measures cancel, and the reaction for the offsets becomes
$F(y,V(y)+a_i(y))$. It is locally bounded: it vanishes near every
nuclear point, and $V$ is smooth off $\operatorname{supp}\nu$.
Apply the first assertion to
the offsets and then add $V$ back.
\end{proof}

